\chapter{Теоретические сведения}
\label{ch:theory}

В этой главе собраны материалы, которые опираются на формулировку учебного задания: сценарий использования модели в приложении, способ кодирования текстовых и музыкальных данных, а также обзор семейств нейросетей, применявшихся в экспериментах~\cite{nikolenko2018}. Формулы даются в минимально необходимом объёме; подробные определения метрик обучения и проверки вынесены в главу~\ref{ch:training}.

\section{Сценарий использования}

Пользователь читает текст и параллельно слушает музыку. Для авторской программы Lyrata предполагается, что выбор книги или главы и выбор подборки треков остаются за человеком; модель может лишь предложить упорядоченный список кандидатов или числовой ``отпечаток'' того, каким должно быть звуковое сопровождение для текущего абзаца. Такой режим мягче полностью автоматического ``ди-джея'', но уже снимает часть рутины: не нужно просматривать весь каталог вручную при каждом переходе между сценами.

С технической стороны в работе фиксируется упрощённая постановка: для каждой строки таблицы данных известен текстовый фрагмент (или его числовое представление) и заранее сопоставленный музыкальный фрагмент, описанный вектором признаков. Модель учится предсказывать этот вектор по тексту. На этапе оценки дополнительно проверяется, попадает ли \textit{фактический} трек в верхнюю долю списка альтернативных треков той же игры при ранжировании по близости предсказанного вектора --- это ближе к тому, как решение может использоваться в интерфейсе.

\section{Игры на Ren'Py как источник пар ``текст--музыка''}

Ren'Py~\cite{srcrenpydocs} --- популярный движок для визуальных новелл. Сценарий игры задаётся скриптом: строки диалогов идут в явном порядке, рядом же автор прописывает команды включения и выключения музыки, звуковых эффектов и переходов между треками. То есть временная последовательность текста и управление звуковым слоем описаны в одном файловом формате, доступном для программного разбора.

Программный парсер может пройти по скрипту, выделить блоки текста между сменами музыкальных состояний и привязать к каждому блоку идентификатор звукового файла или метки из таблицы признаков. Полностью без участия человека обойтись всё равно не удаётся: возможны повторное использование одного трека в разных местах, обрывы или объединение строк, ошибки разметки и различие между ``официальным'' сценарием и тем, как текст воспринимается при чтении. Поэтому датасет дополнительно проходит ручную проверку выборочно или точечную правку правил парсинга; тем не менее доля ручной разметки ``с нуля'' заметно ниже, чем если бы каждую пару приходилось вводить только из аудио и текста без привязки к скрипту движка.

\section{Представление текстовых данных}

Чтобы подавать текст в нейросеть фиксированной входной размерности, каждый фрагмент кодируется вещественным вектором фиксированной длины. В работе используются готовые эмбеддинги русскоязычной модели семейства BERT (RuBERT)~\cite{kuratov2019}: базовая архитектура предобучения ``двунаправленного трансформера'' описана в~\cite{devlin2018}; русскоязычная адаптация используется как готовый генератор признаков длины $768$. Такой способ заменяет ручной отбор признаков (например, частот слов): смысловые закономерности частично уже ``зашиты'' в эмбеддинг на этапе предобучения.

В экспериментах берётся один заранее посчитанный столбец с агрегированием по окну символов вокруг текущей строки (имя столбца, например, \path{text_emb_concat6_stride4_overlap_mean_max512}, в таблице \path{text_musics_rubert}). Это позволяет не пересчитывать тяжёлую языковую модель при каждом запуске обучения регрессии. Если архитектура использует историю предыдущих реплик одной игры, во вход последовательно подаются или объединяются несколько таких векторов подряд; тогда размер входа растёт вместе с числом лагов или обрабатывается рекуррентным блоком.

\section{Представление музыкальных данных}

Каждый используемый в игре трек описывается вектором $y \in \mathbb{R}^{d_y}$ фиксированной размерности, сохранённым в таблице \path{music_data}. Компоненты вектора можно условно разделить на две группы без углубления в название каждой координаты:

\begin{itemize}
	\item \textit{Акустические признаки} --- статистики спектра, мел-кепстральные коэффициенты (MFCC), хрома и родственные величины, полученные инструментами анализа аудио в духе библиотеки librosa~\cite{librosa}. Они отражают темп, спектральную яркость, текстуру звука и устойчивы к конкретному названию файла.
	\item \textit{Теги и метки жанра/настроения} --- числовые столбцы, связанные с внешними разметками (в данных фигурирует набор Jamendo-подобных осей). Они дополняют ``сырой'' звук описанием содержания с точки зрения слушателя.
\end{itemize}

Режим \texttt{-{}-y-group all} объединяет обе группы; в сохранённых прогонах $d_y = 320$. Сужение до только акустики или только тегов возможно параметром пайплайна, но для экспериментов курсовой использовался полный вектор, чтобы сохранить максимум информации для последующего подбора трека.

\section{Архитектуры моделей}

Все рассматриваемые модели решают одну задачу (постановка регрессии с учителем, см.~\cite{nikolenko2018}): по входу из текстовых признаков предсказать вектор $y$. Отличия --- в том, учитывается ли только текущая строка или короткая история прошлых строк этой же игры, и какой блок используется для обработки последовательности. Отдельные учебные конфигурации реализованы скриптами в каталоге \path{text_scripts/ml_scripts/}~\appa{4} общего репозитория проекта~\appa{1}.

\textbf{Линейная модель и многослойный перцептрон.} Базовый уровень --- обучение линейного отображения из $\mathbb{R}^{d_x}$ в $\mathbb{R}^{d_y}$ с регуляризацией (файл \path{train_rubert_to_music_linear.py}, подбор линейного слоя выполняют средствами \texttt{scikit-learn}). Он задаёт нижнюю границу по сложности и скорости. Многослойный перцептрон (MLP, \path{train_rubert_to_music_mlp.py}) строит несколько слоёв вида ``аффинное преобразование + нелинейность'' поверх одного эмбеддинга текущей строки. Формально один скрытый слой можно записать как $h = \sigma(W x + b)$, затем выход $\hat y = V h + c$, где $\sigma$ --- например, функция ReLU. Такая сеть уже может смешивать координаты эмбеддинга нелинейно, но не видит порядка реплик во времени, если не расширить вход.

\textbf{MLP по нескольким лагам эмбеддинга.} Скрипт \path{train_rubert_to_music_mlp_lags.py} объединяет в один длинный вектор эмбеддинги последних $L$ строк текущей игры и подаёт их в MLP. Так модель получает простейший контекст ``что было недавно сказано'', не имея отдельной памяти состояния: вся эта информация выражена одним составным входным вектором --- несколько эмбеддингов подряд записаны в один длинный вектор признаков.

\textbf{Авторегрессия по прошлым значениям цели.} Скрипт \path{train_rubert_to_music_mlp_autoreg_lags.py} добавляет к признакам прошлые векторы $y$ (или их предсказания) для предыдущих шагов. Обучение постепенно переводится с подстановки истинных прошлых значений на всё более частое использование предсказанных, чтобы на длинной последовательности поведение модели было ближе к её работе после обучения, когда истинные прошлые значения уже недоступны.

\textbf{Рекуррентные сети LSTM и GRU.} Файлы \path{train_rubert_to_music_lstm.py} и \path{train_rubert_to_music_gru.py}: последовательность эмбеддингов одной игры обрабатывается рекуррентными ячейками: на каждом шаге обновляется скрытое состояние, ``помнящее'' предыдущий контекст. После последнего шага состояние или выход последней ячейки преобразуются полносвязными слоями в $\hat y$. Такой подход естественен для упорядоченных диалогов.

\textbf{Временная свёрточная сеть (TCN).} Скрипт \path{train_rubert_to_music_tcn.py} заменяет рекуррентность стопкой одномерных свёрток с увеличивающимся шагом (дилатация): каждый слой захватывает всё более широкий временной интервал без явной ячейки памяти. По роли это промежуточный вариант между полносвязной моделью, где на вход подаётся один длинный вектор, собранный из нескольких лагов подряд, и рекуррентной сетью с явным состоянием по времени.

\textbf{Трансформер по ограниченному окну.} Скрипт \path{train_rubert_to_music_transformer_window.py} обрабатывает последние позиции через механизм самовнимания~\cite{vaswani2017} внутри окна ограниченной длины (параметр \path{tf_window}). Так модель может выборочно усиливать вклад отдалённых по времени реплик по сравнению с фиксированной свёрткой, но цена --- больший число параметров и риск переобучения на конкретные игры~\cite{nikolenko2018}.

Численные гиперпараметры (число слоёв, ширина, число лагов) для каждого сохранённого прогона указаны в файле \texttt{config.json} каталога прогона~\appa{5} и в манифесте \texttt{manifest.jsonl}~\appa{3}. Итоговое сравнение по качеству приводится в главе~\ref{ch:training}.
